%!TEX root = ../thesis.tex
% FuzzerGym for \cref{sec:rw-mdp-formulations}: redrawn after the architecture
% overview of Drozd and Wagner, "FuzzerGym" (arXiv:1807.07490).
% Layout rule: the fuzzing loop is a rectangular cycle traversed
% counter-clockwise --- up the left column, right along the top, down the right
% column and back left along the bottom --- so that every edge stays horizontal
% or vertical. The two arrows between the panels are what fixes the vertical
% placement of everything else: the state arrow runs at y=0, the height of both
% the environment box and the instrumented target, and the action arrow at
% y=-4.1, the height of both the Q-value output and the mutator selection. Move
% one of those four nodes and an arrow goes diagonal.
% Spacing constraint: an edge label sits in the 1.6cm gap between the two
% columns of the loop, so it must be narrower than that gap or it slides under
% the boxes; "select a random input" is stacked over two lines for that reason.
% The whole picture is 14.3cm wide against a 14.56cm text width, so there is no
% room to widen a column without taking the width back somewhere else.
% Deviations from the original: its promotional bullet lists ("competitive
% fuzzing via OpenAI Gym", "leverage numerous RL implementations", ...) are
% dropped, its separate "corpus input" box is folded into the edge label
% "select a random input", and the layer stack is drawn as boxes rather than as
% unit graphs.
\begin{figure}[htbp]
  \centering
  \begin{tikzpicture}[
      >={Stealth[round]}, font=\footnotesize,
      panel/.style  = {draw=black!45, densely dashed, rounded corners=3pt,
                       inner sep=3.5mm},
      ptitle/.style = {font=\footnotesize\bfseries},
      agent/.style  = {draw, rounded corners=2pt, fill=blue!8, align=left,
                       text width=39mm, font=\scriptsize, inner sep=2mm},
      layer/.style  = {draw, rounded corners=1pt, fill=blue!8, align=center,
                       font=\scriptsize, minimum height=4.5mm,
                       minimum width=30mm, inner sep=1pt},
      netbox/.style = {draw=black!45, rounded corners=2pt, inner sep=1.5mm},
      sig/.style    = {align=center, font=\scriptsize, inner sep=1.5pt},
      proc/.style   = {draw, fill=white, align=center, text width=17mm,
                       minimum height=9mm, font=\scriptsize, inner sep=1.5mm},
      store/.style  = {draw, fill=black!8, align=center, text width=17mm,
                       minimum height=9mm, font=\scriptsize, inner sep=1.5mm},
      act/.style    = {draw, fill=blue!8, align=center, text width=17mm,
                       minimum height=9mm, font=\scriptsize, inner sep=1.5mm},
      dec/.style    = {draw, ellipse, fill=black!8, align=center,
                       font=\scriptsize, inner sep=1pt},
      bridge/.style = {draw, rounded corners=2pt, fill=black!5, align=center,
                       font=\scriptsize, text width=20mm, inner sep=1.5mm},
      flow/.style   = {->, thick, black!70},
      link/.style   = {->, very thick, KITgreen},
      lbl/.style    = {font=\scriptsize, inner sep=2.5pt, align=center},
    ]
    % --- the agent: one Gym environment and the network behind it ------------
    \node[ptitle] (gymtitle) at (2.15, 1.75) {FuzzerGym};
    \node[agent]  (env)      at (2.15, 0)    {%
      \frenchspacing\raggedright  % no extra space after the colons, at 7pt it reads as a gap
      \textbf{OpenAI Gym environment}
      \begin{itemize}[label=\textbullet, nosep, topsep=2pt,
                      leftmargin=1.1em, labelsep=0.3em]
        \item state: the test input as a bit array
        \item reward: newly covered lines
        \item action: a libFuzzer mutation operator
      \end{itemize}};

    \node[sig]   (qtitle) at (2.15,-1.65) {deep double $Q$-network};
    \node[layer] (l1)     at (2.15,-2.15) {LSTM, $64\times 64$};
    \node[layer] (l2)     at (2.15,-2.70) {MLP, $64\times 64$};
    \node[layer] (l3)     at (2.15,-3.25) {MLP, $64\times 64$};
    \node[netbox, fit=(qtitle)(l1)(l3)] (qnet) {};
    \node[sig]   (q)      at (2.15,-4.10) {a $Q$-value per operator};

    \draw[flow] (env)  -- (qnet);
    \draw[flow] (qnet) -- (q);

    % --- the fuzzer: libFuzzer's own loop, one step of it given away --------
    \node[ptitle] (libtitle) at (10.6, 1.75) {libFuzzer};
    \node[proc]   (code)     at ( 8.8, 0)    {instrumented\\target};
    \node[dec]    (covq)     at (12.4, 0)    {new\\coverage?};
    \node[store]  (corpus)   at (12.4,-2.05) {input\\corpus};
    \node[proc]   (test)     at ( 8.8,-2.05) {test input};
    \node[act]    (mut)      at ( 8.8,-4.10) {select a\\mutator};

    \draw[flow] (mut)  -- node[lbl,right] {mutate the input}   (test);
    \draw[flow] (test) -- node[lbl,right] {run the test}       (code);
    \draw[flow] (code) --                                      (covq);
    \draw[flow] (covq) -- node[lbl,left]  {add to the\\corpus} (corpus);
    \draw[flow,rounded corners=3pt] (corpus.south) -- (12.4,-4.10)
                       -- node[lbl,above] {select a\\random input} (mut.east);

    % --- the two panels, drawn once every node inside them exists ------------
    \coordinate (gymfoot) at (2.15,-4.55);   % keeps both panels the same height
    \node[panel, fit=(gymtitle)(env)(qnet)(q)(gymfoot)]              (gym) {};
    \node[panel, fit=(libtitle)(code)(covq)(corpus)(test)(mut)]      (lib) {};

    % --- the bridge, and the only two messages that cross it -----------------
    \node[bridge] (grpc) at (6.05,-2.05) {gRPC bridge:\\asynchronous\\streaming};
    \draw[dotted,black!55] (grpc.north) -- (6.05, 0);
    \draw[dotted,black!55] (grpc.south) -- (6.05,-4.10);

    \draw[link] (lib.west |- code) --
                node[lbl,above] {state and reward}          (gym.east |- env);
    \draw[link] (gym.east |- q)    --
                node[lbl,below] {action:\\a mutation operator} (lib.west |- mut);
  \end{tikzpicture}
  \caption[The FuzzerGym architecture]{%
    FuzzerGym, redrawn after \textcite{drozd_fuzzergym_2018}.
    libFuzzer keeps its own fuzzing loop, and the single step it hands over ---
    the choice of the mutation operator --- becomes the action of an OpenAI Gym
    environment.
    The agent observes the test input as a bit array together with the program
    state the LLVM sanitizers expose, is rewarded with the number of newly
    covered lines, and scores the thirteen operators with a deep double
    $Q$-network whose first layer is recurrent.
    Both directions cross the language boundary through an asynchronous gRPC
    bridge, so that libFuzzer never has to block on a forward pass.
  }
  \label{fig:rw-fuzzergym}
\end{figure}
